\newpage

\section{Заключение}
%\begin{table}[h!t]
%\begin{center}
%\caption{Результаты работы алгоритма}
%\label{table_2}
%\begin{tabular}{|c|c|c|c|c|}
%\hline
%	Ряд,~$\textbf{x}$ &Длина ряда,~$N$& \# сегментов,~$K$&Длина сегмента,~$T$& Ошибка,~$S$\\
%	\hline
%	\multicolumn{1}{|l|}{Phys.~Motion~1}
%	& 900& 2& 40& 0.06\\
%	\hline
%	\multicolumn{1}{|l|}{Phys.~Motion~2}
%	& 900& 2& 40& 0.03\\
%	\hline
%	\multicolumn{1}{|l|}{Synthetic~1}
%	& 2000& 2& 20& 0.04\\
%	\hline
%	\multicolumn{1}{|l|}{Synthetic~2}
%	& 2000& 3& 20& 0.03\\
%\hline

%\end{tabular}
%\end{center}
%\end{table}


В работе рассматривалась задача прореживания моделей нейросетей. Рассматривался метод оптимального прореживания и метод, основаный на вариационном подходе. Был предложен алгоритм прореживания, основаный на методе Белсли для удаления зависимых параметров модели. В ходе эксперимента было показано, что нейросети прорежены методом Белсли являются более устойчивы к шуму на входных данных. Качества прогноза нейросетей после прореживания методом Белсли не хуже качества прогноза нейросетей, полученных другими методами.
